% Card: the long side (rook and pawn vs rook)

Every pawn except a rook's pawn divides the board into two parts of different width: the
\vocab{short-side}{short side}{The narrower of the two parts into which a pawn (other than a
rook's pawn) divides the board} and the
\vocab{long-side}{long side}{The wider of the two parts into which a pawn (other than a rook's
pawn) divides the board. The defending rook belongs there, so that it can give side checks}.
In \xref{diagram}{side-checks} the e-pawn has three files to its right (f, g, h) and four to its
left (a, b, c, d).

Because \vocab{side-checks}{side checks}{} need room, the defending rook belongs on the long
side. The defending king goes to the short side, where it does not stand in the way of its own
rook. Put Black's rook on h2 instead of a2 in \xref{diagram}{side-checks}: now it is on the short
side, too close to the pawn to give side checks, and White wins.
